%% style: lecture
%% desc: 讲义/课堂笔记风 —— 面向学习：目标先行、分节讲解、重点框、练习提示
%% 用法：AGENT 读此模板，以 MD 底稿为信息源、按此结构生成 note.tex（不直接替换占位符）。
%% 占位符 %%TITLE%% / %%CONTENT%% 标出内容插入点。
\documentclass[11pt]{article}
\usepackage[UTF8, scheme=plain]{ctex}
\usepackage[margin=2.5cm]{geometry}
\usepackage{amsmath, amssymb}
\usepackage{enumitem}
\usepackage{tcolorbox}
\usepackage{hyperref}
\hypersetup{hidelinks}

\title{\textbf{讲义：%%TITLE%%}}
\author{AI 笔记}
\date{\today}

\newtcolorbox{keypoint}{colback=blue!5!white, colframe=blue!60!black, title=重点}
\newtcolorbox{note}{colback=yellow!10!white, colframe=orange!70!black, title=注意}

\begin{document}
\maketitle

\section{学习目标}
%% 用 3-5 条列出本讲完成后应掌握的内容

\section{背景与引入}
%% 为什么讲这个：问题/场景/动机（从底稿开头提炼）

\section{主体内容}
%% 按底稿顺序分节 \subsection{}，逐步展开：
% - 概念先定义（\emph{}）
% - 再给例子（\begin{example} 或列表）
% - 关键点放 keypoint 框

\begin{keypoint}
%% 每节的要点浓缩
\end{keypoint}

\begin{note}
%% 易错点/陷阱/需注意处
\end{note}

\section{小结}
%% 用「学了什么」清单式总结

\section{思考练习}
%% 1-3 道基于内容的思考题（从底稿可回答）

\end{document}
